\documentclass[11pt]{article}

\usepackage[margin=1in]{geometry}
\usepackage{amsmath}
\usepackage{amssymb}
\usepackage{booktabs}
\usepackage{array}
\usepackage{enumitem}
\usepackage[colorlinks=true,linkcolor=blue,urlcolor=blue]{hyperref}

\setlength{\parindent}{0pt}
\setlength{\parskip}{0.55em}
\setlist[enumerate]{leftmargin=2.2em,itemsep=0.35em,topsep=0.25em}

\title{S181M116 Derivatives and Alternative Investments\\
Exam Solution}
\author{Swapnil Singh}
\date{}

\begin{document}
\maketitle

This answer key applies to both the English and Lithuanian multiple-choice
versions of Variant B. Each item is worth 4 points.

\section*{Answer Summary}

\begin{center}
\begin{tabular}{cccccccccc}
\toprule
1.1 & 1.2 & 1.3 & 1.4 & 2.1 & 2.2 & 2.3 & 2.4 & 3.1 & 3.2 \\
\midrule
B & A & B & A & A & C & B & B & A & C \\
\bottomrule
\end{tabular}

\vspace{0.75em}

\begin{tabular}{cccccccccc}
\toprule
3.3 & 3.4 & 4.1 & 4.2 & 4.3 & 4.4 & 5.1 & 5.2 & 5.3 & 5.4 \\
\midrule
B & A & D & B & A & A & C & A & B & C \\
\bottomrule
\end{tabular}
\end{center}

\section*{Question 1: Concepts, Payoffs, and Market Structure}

\begin{enumerate}[label=\textbf{1.\arabic*.}]
    \item \textbf{B.} The payoff $\max(90-S_T,0)$ depends on the future asset
    price $S_T$, so the contract is a derivative. The payoff is put-like, not
    linear like a futures payoff.

    \item \textbf{A.} The call payoff is
    \[
        \max(88-75,0)=13,
    \]
    so the call profit is $13-6=7$. The put payoff is
    \[
        \max(75-88,0)=0,
    \]
    so the put profit is $0-4=-4$.

    \item \textbf{B.} OTC contracts can be customized to the counterparties'
    needs, but this customization usually brings more counterparty, valuation,
    and liquidity risk than exchange-traded contracts.

    \item \textbf{A.} Futures are marked to market. If the futures position loses
    money, cash must be posted through variation margin even if the underlying
    exposure will generate an offsetting gain only later.
\end{enumerate}

\section*{Question 2: Futures Settlement and Forward Pricing}

\begin{enumerate}[label=\textbf{2.\arabic*.}]
    \item \textbf{A.} A short futures position gains when the futures price
    falls:
    \[
        8 \times 25{,}000 \times (4.12-4.05)
        = 200{,}000 \times 0.07
        = 14{,}000.
    \]
    The result is a USD 14,000 gain.

    \item \textbf{C.} The margin account receives the USD 14,000 gain:
    \[
        38{,}000+14{,}000=52{,}000.
    \]
    This is above the USD 30,000 maintenance margin, so there is no margin call.

    \item \textbf{B.} For a stock with known cash dividends,
    \[
        F_0=(S_0-I)e^{rT}
        =(120-2.40)e^{0.035\times 0.5}
        =119.68.
    \]

    \item \textbf{B.} The quoted forward price 121 is above the no-arbitrage
    forward price 119.68. The cash-and-carry arbitrage is to buy the asset,
    short the forward, and lock in about
    \[
        121-119.68=1.32
    \]
    per share.
\end{enumerate}

\section*{Question 3: Cross Hedging, Basis Risk, and Effective Prices}

\begin{enumerate}[label=\textbf{3.\arabic*.}]
    \item \textbf{A.} The manufacturer will buy copper later, so it is exposed to
    rising copper prices. A long futures position gains when the futures price
    rises and therefore hedges that risk.

    \item \textbf{C.} The minimum-variance hedge ratio is
    \[
        h^*=\rho\frac{\sigma_S}{\sigma_F}
        =0.72\frac{0.18}{0.16}
        =0.8100.
    \]

    \item \textbf{B.} The number of contracts before rounding is
    \[
        N^*=h^*\frac{Q_A}{Q_F}
        =0.81\frac{1{,}250{,}000}{25{,}000}
        =0.81\times 50
        =40.5.
    \]

    \item \textbf{A.} For the long futures hedge, the futures gain per pound is
    \[
        F_2-F_1=4.05-3.80=0.25.
    \]
    The effective purchase price is therefore
    \[
        S_2-(F_2-F_1)=4.12-0.25=3.87.
    \]
    The realized basis is
    \[
        S_2-F_2=4.12-4.05=0.07.
    \]
    Since $0.07>0.03$, the realized basis is worse for the buyer than expected.
\end{enumerate}

\section*{Question 4: Option No-Arbitrage and Binomial Pricing}

\begin{enumerate}[label=\textbf{4.\arabic*.}]
    \item \textbf{D.} Put-call parity for a non-dividend-paying stock is
    \[
        c+Ke^{-rT}=p+S_0.
    \]
    Thus
    \[
        p=c+Ke^{-rT}-S_0
        =2.60+50e^{-0.05\times0.25}-48
        =3.98.
    \]

    \item \textbf{B.} The protective put costs
    \[
        S_0+p=48+3.20=51.20.
    \]
    The fiduciary call costs
    \[
        c+Ke^{-rT}=2.60+50e^{-0.0125}=51.98.
    \]
    These two portfolios have the same payoff, but the fiduciary call is more
    expensive. Sell the fiduciary call and buy the protective put.

    \item \textbf{A.} The one-step stock prices are
    \[
        S_u=48(1.18)=56.64,\qquad S_d=48(0.88)=42.24.
    \]
    The call payoffs are
    \[
        C_u=\max(56.64-50,0)=6.64,\qquad
        C_d=\max(42.24-50,0)=0.
    \]

    \item \textbf{A.} The risk-neutral probability is
    \[
        q=\frac{e^{rT}-d}{u-d}
        =\frac{e^{0.0125}-0.88}{1.18-0.88}
        =0.4419.
    \]
    The call value is
    \[
        e^{-rT}\left(qC_u+(1-q)C_d\right)
        =e^{-0.0125}(0.4419\times6.64)
        =2.90.
    \]
\end{enumerate}

\section*{Question 5: Option Bounds, Early Exercise, and Strategy Logic}

\begin{enumerate}[label=\textbf{5.\arabic*.}]
    \item \textbf{C.} For a European call on a non-dividend-paying stock,
    \[
        c\geq \max(S_0-Ke^{-rT},0).
    \]
    Here
    \[
        Ke^{-rT}=68e^{-0.045\times0.75}=65.74,
    \]
    so the lower bound is
    \[
        70-65.74=4.26.
    \]
    A quote of $c=3.80$ violates the lower bound.

    \item \textbf{A.} For a European put,
    \[
        p\geq \max(Ke^{-rT}-S_0,0)
        =\max(65.74-70,0)
        =0.
    \]
    A quote of $p=0.60$ does not violate this lower bound.

    \item \textbf{B.} Put-call parity gives
    \[
        p=c+Ke^{-rT}-S_0
        =6.00+65.74-70
        =1.74.
    \]

    \item \textbf{C.} An American put must be worth at least as much as the
    otherwise identical European put because it has all the European exercise
    rights plus the option to exercise early. A covered call is not risk-free:
    owning the stock and writing a call still leaves downside stock risk.
\end{enumerate}

\end{document}
